% !TEX root = ../main.tex
\chapter{Referência da biblioteca padrão}
\label{ap:stdlib}

Todas as funções intrínsecas da \cmm{}, agrupadas. \enquote{Tipos} indica os tipos aceitos; erros de tipo são apontados na compilação.

\section{Entrada e saída}

\begin{longtable}{@{}llp{0.52\linewidth}@{}}
\toprule
\textbf{Função} & \textbf{Tipos} & \textbf{Descrição} \\
\midrule
\endhead
\code{in(p)} & \code{int} & Lê a porta de entrada \code{p}. \\
\code{fin(p)} & \code{float} & Lê a porta \code{p} convertendo para \emph{float}. \\
\code{out(p, x);} & qualquer & Escreve \code{x} na porta de saída \code{p}. \\
\code{fout(p, x);} & \code{float} & Escreve em formato \emph{float}. \\
\bottomrule
\end{longtable}

\section{Funções especiais (economia de hardware)}

\begin{longtable}{@{}llp{0.52\linewidth}@{}}
\toprule
\textbf{Função} & \textbf{Tipos} & \textbf{Descrição} \\
\midrule
\endhead
\code{norm(x)} & \code{int} & $x\,/\,$\texttt{\#NUGAIN} sem instanciar divisor. \\
\code{pset(x)} & \code{int}, \code{float} & $x$ se $x \geq 0$; senão $0$. \\
\code{abs(x)} & todos & Valor absoluto; magnitude para \code{comp}. \\
\code{sign(x, y)} & \code{int}, \code{float} & $y$ com o sinal de $x$. \\
\code{copy(x, y);} & todos & Cópia bit a bit, sem checagem de tipo. \\
\bottomrule
\end{longtable}

\section{Não-lineares (macros de assembly)}

\begin{longtable}{@{}lp{0.64\linewidth}@{}}
\toprule
\textbf{Função} & \textbf{Descrição} \\
\midrule
\endhead
\code{sqrt(x)} & Raiz quadrada (Newton--Raphson). Não aceita \code{comp}. \\
\code{sin(x)}, \code{cos(x)}, \code{tan(x)} & Trigonométricas. \\
\code{atan(x)} & Arco-tangente. \\
\code{sinh(x)}, \code{cosh(x)}, \code{tanh(x)} & Hiperbólicas. \\
\code{exp(x)}, \code{log(x)} & Exponencial e logaritmo natural. Não aceitam \code{comp}. \\
\code{pow(x, y)} & Potência; estratégia depende do expoente (constante inteiro $\to$ multiplicações; variável inteiro $\to$ laço; senão $e^{y\ln x}$). \\
\code{floor(x)}, \code{ceil(x)}, \code{round(x)} & Arredondamentos (retornam \code{float}; \code{round} afasta do zero no meio). \\
\bottomrule
\end{longtable}

\section{Complexos}

\begin{longtable}{@{}lp{0.64\linewidth}@{}}
\toprule
\textbf{Função} & \textbf{Descrição} \\
\midrule
\endhead
\code{real(z)}, \code{imag(z)} & Partes real e imaginária. \\
\code{fase(z)} & Argumento (ângulo). \\
\code{mod2(z)} & Magnitude ao quadrado. \\
\code{complex(re, im)} & Constrói um \code{comp}. \\
\code{conj(z)} & Conjugado. \\
\bottomrule
\end{longtable}

\section{Notação de Dirac}

\begin{longtable}{@{}lp{0.60\linewidth}@{}}
\toprule
\textbf{Sintaxe} & \textbf{Operação} \\
\midrule
\endhead
\texttt{<a|b>} & Produto interno (expressão). \\
\texttt{a \# |0>;} & Zera o vetor. \\
\texttt{a \# |M|b>;} & Matriz vezes vetor. \\
\texttt{a \# c|b>;} & Vetor escalado. \\
\texttt{A \# |a><b|;} & Produto externo. \\
\texttt{out(p, c|a>);} & Emite o vetor escalado pela porta \code{p}. \\
\bottomrule
\end{longtable}

\section{Famílias de instruções do processador}

Para leitura das trilhas de assembly nas ondas: os mnemônicos seguem convenções de prefixo/sufixo. \code{F\_} indica a versão em ponto flutuante; \code{S\_}/\code{SF\_}, operando vindo da pilha; \code{P\_}, empilha antes; \code{\_M}, opera sobre memória; \code{\_V}, variante vetorial/indexada.

\begin{longtable}{@{}lp{0.62\linewidth}@{}}
\toprule
\textbf{Família} & \textbf{Mnemônicos principais} \\
\midrule
\endhead
Carga/armazenamento & \code{LOD}, \code{SET}, \code{LDI}/\code{STI} (indexados), \code{ILI}/\code{ISI} (bit-reverso) \\
Pilha & \code{PSH}, \code{POP} \\
E/S & \code{INN}, \code{F\_INN}, \code{OUT} \\
Fluxo & \code{JMP}, \code{JIZ}, \code{CAL}, \code{RET}, \code{NOP} \\
Aritmética inteira & \code{ADD}, \code{MLT}, \code{DIV}, \code{MOD}, \code{NEG}, \code{ABS}, \code{SGN}, \code{PST}, \code{NRM} \\
Aritmética float & \code{F\_ADD}, \code{F\_SU1}/\code{F\_SU2}, \code{F\_MLT}, \code{F\_DIV}, \code{F\_NEG}, \code{F\_ABS}, \code{F\_ROT}, \code{F\_SCL}, \code{XPO} \\
Conversões & \code{I2F}, \code{F2I} \\
Lógica e bits & \code{AND}, \code{ORR}, \code{XOR}, \code{INV}, \code{LAN}, \code{LOR}, \code{LIN} \\
Comparações & \code{LES}, \code{GRE}, \code{EQU} (e as versões \code{F\_}) \\
Deslocamentos & \code{SHL}, \code{SHR}, \code{SRS} (equivale ao operador \code{>>>}) \\
\bottomrule
\end{longtable}

Cada mnemônico usado pela primeira vez no programa liga o bloco de hardware correspondente no processador gerado --- a lista completa e o percentual de uso da ISA aparecem no terminal TASM a cada compilação.
